\begin{ouverture}[grandeurs]{12}{Masses et contenances}{Partie 3 · Grandeurs et mesures}
{Grammes · Kilogrammes · Litres · Centilitres}
\objectif{Choisir une unité de masse ou de contenance adaptée.}
\objectif{Convertir en s’appuyant sur les relations entre unités.}
\objectif{Résoudre des problèmes de masses et de contenances.}
\prerequis{Multiplication par 10, 100, 1 000 et décimaux.}
\end{ouverture}
\pagestyle{otMain}\essentiel\fontsize{11}{15}\selectfont
\phantomsection\label{cours-12-1}
\essSection{1}{Les masses}
\begin{definition}Les unités de masse vont du milligramme à la tonne : mg, g, kg, t. Relations à connaître : $1\ \mathrm g=1\,000\ \mathrm{mg}$, $1\ \mathrm{kg}=1\,000\ \mathrm g$ et $1\ \mathrm t=1\,000\ \mathrm{kg}$.\end{definition}
\begin{exemple}Références utiles : une pièce de 1 euro pèse environ $7{,}5$ g, une pomme environ 150 g, une bouteille d’eau pleine environ $1{,}5$ kg, une petite voiture environ $1{,}2$ t.\end{exemple}
\begin{savoirfaire}[Convertir des masses]$2{,}5$ kg $=2\,500$ g, car 1 kg vaut $1\,000$ g et $2{,}5\times1\,000=2\,500$. Dans l’autre sens, 500 g est la moitié de $1\,000$ g : c’est $0{,}5$ kg.\end{savoirfaire}
\begin{attention}Sur une balance à plateaux, l’équilibre compare deux masses : l’objet pèse autant que l’ensemble des masses marquées posées sur l’autre plateau.\end{attention}
\clearpage
\phantomsection\label{cours-12-2}
\essSection{2}{Les contenances}
\begin{definition}Les unités de contenance vont du millilitre à l’hectolitre : mL, cL, dL, L, hL. Relations à connaître : $1\ \mathrm L=100\ \mathrm{cL}=1\,000\ \mathrm{mL}$ et $1\ \mathrm{hL}=100\ \mathrm L$.\end{definition}
\begin{exemple}Références utiles : une cuillère à café contient environ 5 mL, un verre environ 20 cL, une grande bouteille $1{,}5$ L, un seau environ 10 L, une baignoire environ 150 L.\end{exemple}
\begin{savoirfaire}[Convertir des contenances]$2$ L $=200$ cL, car 1 L vaut 100 cL. Et 350 cL $=3{,}50$ L, car 100 cL font 1 L : on cherche combien de fois 100 tient dans 350.\end{savoirfaire}
\begin{attention}Comparer des contenances exige la même unité : 75 cL et 1 L se comparent en centilitres, $75<100$. La forme du récipient ne dit pas sa contenance.\end{attention}

\clearpage\rubrique{Découverte}
\begin{blocLivre}{ex-12-1}\begin{exercice}[1]{La bonne unité}
Choisir l’unité adaptée (mg, g, kg, t ou mL, cL, L) : masse d’une gomme ; masse d’un camion chargé ; contenance d’une cuillère ; contenance d’une piscine.
\end{exercice}\end{blocLivre}
\begin{blocLivre}{ex-12-2}\begin{exercice}[1]{Convertir des masses}
Convertir : 3 kg en g ; $2{,}5$ kg en g ; 500 g en kg.
\end{exercice}\end{blocLivre}
\begin{blocLivre}{ex-12-3}\begin{exercice}[1]{Convertir des contenances}
Convertir : 2 L en cL ; 350 cL en L ; 5 hL en L.
\end{exercice}\end{blocLivre}
\begin{blocLivre}{ex-12-4}\begin{exercice}[1]{Comparer}
Comparer $1{,}2$ kg et 950 g, puis 75 cL et 1 L.
\end{exercice}\end{blocLivre}
\begin{blocLivre}{ex-12-5}\begin{exercice}[1]{Des estimations sensées}
Associer chaque objet à une mesure : une orange ; un cartable rempli ; une canette de soda ; une baignoire. Mesures proposées : 33 cL ; 150 L ; 200 g ; 7 kg.
\end{exercice}\end{blocLivre}

\rubrique{Entraînement}
\begin{blocLivre}{ex-12-6}\begin{exercice}[2]{La balance à plateaux}
Une balance est en équilibre : à gauche, un sac de farine ; à droite, des masses de 500 g, 200 g et 50 g. Quelle est la masse du sac ?
\label{fig-ex-12-6}\figureGeo{mas-6}
\end{exercice}\end{blocLivre}
\begin{blocLivre}{ex-12-7}\begin{exercice}[2]{Le colis complet}
Un colis contient un livre de $1{,}2$ kg et un jeu de 850 g. Calculer la masse totale, en grammes puis en kilogrammes.
\end{exercice}\end{blocLivre}
\begin{blocLivre}{ex-12-8}\begin{exercice}[2]{Des verres bien remplis}
Une carafe contient $1{,}5$ L de jus. Combien de verres de 25 cL peut-on remplir ?
\end{exercice}\end{blocLivre}
\begin{blocLivre}{ex-12-9}\begin{exercice}[2]{Le chargement du camion}
Un camion peut charger 1 t. On pose une palette de 850 kg puis une caisse de 325 kg. Le chargement est-il possible ?
\end{exercice}\end{blocLivre}
\begin{blocLivre}{ex-12-10}\begin{exercice}[2]{La recette doublée}
Une recette demande 250 g de farine. On double la recette, avec un paquet de 1 kg. Quelle masse de farine utilise-t-on et combien en reste-t-il ?
\end{exercice}\end{blocLivre}

\rubrique{Problèmes}
\begin{blocLivre}{ex-12-11}\begin{exercice}[3]{Les briques de lait}
La cantine reçoit 3 packs de 6 briques de 1 L de lait. Elle utilise 15 L dans la semaine. Quelle contenance de lait reste-t-il ?
\end{exercice}\end{blocLivre}
\begin{blocLivre}{ex-12-12}\begin{exercice}[3]{L’ascenseur}
Un ascenseur accepte 300 kg au plus. Quatre personnes de 72 kg chacune y montent avec un colis de 15 kg. Est-ce autorisé ?
\end{exercice}\end{blocLivre}
\begin{blocLivre}{ex-12-13}\begin{exercice}[3]{La citerne de l’école}
Une citerne de récupération d’eau contient 5 hL. On en tire 120 L pour arroser. Quelle contenance d’eau reste-t-il, en litres ?
\end{exercice}\end{blocLivre}
\begin{blocLivre}{ex-12-14}\begin{exercice}[3]{Les pots de confiture}
Avec $2{,}4$ kg de confiture, combien de pots de 300 g peut-on remplir exactement ?
\end{exercice}\end{blocLivre}
\begin{blocLivre}{ex-12-15}\begin{exercice}[3]{Plumes ou plomb ?}
Sacha demande : « Qu’est-ce qui est le plus lourd : 1 kg de plumes ou 1 kg de plomb ? » Répondre, puis expliquer pourquoi un litre d’eau est une bonne référence de masse.
\end{exercice}\end{blocLivre}

\clearpage
